\documentclass{../../../Styles/MS-Style}
\addbibresource{../../../Assets/Citations/references.bib}

\usepackage[
    pdftitle={André Mossi CV (EN)},
    pdfauthor={André Mossi},
    pdfcreator={LaTeX and NeoVim}
]{hyperref}

\begin{document}

\begin{center}
	\authorTitle{Roberto (André) Mossi Milla}

	\contactBlock{
		Erie, Pennsylvania, 16506
	}{
		814-790-1591
	}{
		mossiroberto0392@gmail.com
	}{
		andremossi-linktree.vercel.app
	}{
		github.com/AndreM222
	}
\end{center}

\begin{Form}
	\sectionTitle{SUMMARY}
    \normalsize{
        Software Engineer experienced in AI-powered full-stack platforms,
        computer vision systems, and real-time simulation environments.
        Background includes production web applications, artificial intelligence research,
        cloud integrations, and end-to-end system development across healthcare,
        industrial automation, and multimodal AI workflows.
    }

	\sectionTitle{WORK EXPERIENCE}

    \subSectionTitle{Ingeniería Digital}{Tegucigalpa, Honduras}

    \subSubSectionTitle{Full-Stack Developer / AI Systems Engineer}{May 2026 - Present}
    \textbf{AI-Assisted Medical Document Generation Platform}
    \begin{listBlock}
        \generalItem
        Developed an AI-assisted medical document generation platform for post-surgery reports (IDR),
        combining OCR extraction pipelines, LLM-powered grammar correction, and automated professional
        document structuring using Next.js, React, TypeScript, PostgreSQL, AWS S3,
        Amazon Textract, and OpenAI APIs.

        \generalItem
        Architected authenticated document management workflows using Auth0, enabling users
        to generate, analyze, edit, version, and manage AI-generated reports through
        a full-stack web platform.

        \generalItem
        Designed scalable file-processing pipelines integrating OCR extraction,
        cloud storage, AWS S3 buckets, and PostgreSQL persistence for intelligent
        document retrieval and processing.

        \generalItem
        Implemented internationalization workflows and automated translation asset generation
        using CMake to support multilingual deployments across US and Honduran clients.
    \end{listBlock}

    \textbf{AI Vehicle Damage Assessment Platform}
    \begin{listBlock}
        \generalItem
        Developed an AI-assisted vehicle damage assessment platform utilizing multimodal AI workflows,
        VIN history analysis, image-based damage detection, and insurance estimation pipelines.

        \generalItem
        Integrated computer vision and LLM APIs to analyze uploaded vehicle imagery,
        validate VIN-related information, compare against historical vehicle records,
        and support insurance decision workflows.

        \generalItem
        Developed data acquisition services using Playwright to collect and process
        vehicle information required for valuation and insurance analysis workflows.

        \generalItem
        Implemented interactive damage annotation and review tools using Konva,
        allowing users and insurance personnel to identify, visualize, and validate
        detected vehicle damage regions.

        \generalItem
        Developed responsive client-facing and administrative dashboards using React,
        Next.js, TypeScript, PostgreSQL, AWS S3, and Shadcn/UI.

        \generalItem
        Maintained engineering standards and CI/CD workflows using GitHub Actions,
        Husky, ESLint, Prettier, lint-staged, and Commitizen.
    \end{listBlock}

	\subSectionTitle{Gannon University}{Pennsylvania, United States}

	\subSubSectionTitle{Computer Vision Engineer (Project-Based) Industry Partner (NDA)}{August 2024 - May 2025}
	\begin{listBlock}
		\generalItem
		Architected and deployed a real-time industrial traceability system using Ultralytics YOLO
		to detect and track pallets, automatically generating structured events with images
		and detection metadata.\cite{ComputerVisionVisual}

		\generalItem
		Collaborated in a cross-functional team contributing across hardware, backend, and software components;
		took primary ownership of the monitoring dashboard and visualization interface.

		\generalItem
		Developed a full-stack web dashboard (React, JavaScript) integrated with backend services and TimescaleDB
		to visualize detection events, images, and operational analytics in real time.

		\generalItem
		Implemented administrative tools within the dashboard to delete images, correct detections,
		and curate labeled datasets used for iterative object detection model retraining,
		improving dataset quality and system reliability.

		\generalItem
		Engineered an event-driven pipeline that synchronized frame capture with detection metadata,
		enabling downstream anomaly and damage analysis workflows.

		\generalItem
		Delivered recurring technical presentations and system demonstrations to industry stakeholders
		and internal reviewers, supporting continued project funding and informing architectural decisions,
		including migration to NVIDIA Jetson hardware for improved performance and thermal stability.

		\generalItem
		Containerized distributed services using Docker and integrated multi-language components
		(C++, Python, Rust) to coordinate inference, data processing, and system communication.
	\end{listBlock}

	\subSubSectionTitle{Research Assistant}{August 2024 - May 2025}
	\begin{listBlock}
		\generalItem
		Conducted research under faculty supervision on multi-agent artificial intelligence
		simulating natural selection and survival behaviors in dynamic environments.\cite{TagResearchRecommendation}

		\generalItem
		Implemented NeuroEvolution of Augmenting Topologies (NEAT) to evolve neural network architectures and
		behavioral strategies across generations.

		\generalItem
		Developed a 3D multi-agent simulation environment in Unreal Engine featuring
		procedural world generation, physics-based interactions, and reinforcement learning pipelines
		using C++, CUDA, and PyTorch.

		\generalItem
		Research paper accepted and published by the American Society for Engineering Education
		(ASEE) and presented at a Sigma Xi conference at Penn State University.\cite{AITagResearchPublication} \cite{AITagSigmaXIPresentation}
	\end{listBlock}

	\subSectionTitle{Dinant}{Tegucigalpa, Honduras}

	\subSubSectionTitle{Full-Stack Developer}{June - July 2023}
	\begin{listBlock}
		\generalItem
		Led the end-to-end design and development of two production web platforms, conducting requirements
		gathering, system planning, database architecture design, client presentations, and final
		deployment.\cite{DinantRecommendation}

		\generalItem
		Architected and implemented a contract lifecycle management system for a device sales, rental, and repair
		company. Designed relational SQL schemas to track device state (sold, rented, under repair, returned),
		service eligibility, country-based filtering, and client records across multiple regions.

		\generalItem
		Developed a multi-role web application (Oracle APEX, JavaScript, HTML, CSS) enabling guest contract signing,
		digital signature capture, automated PDF generation, and administrative dashboards with advanced filtering
		by date, country, device state, and contract status.

		\generalItem
		Engineered a repair workflow system with state transitions, damage logging, cost tracking, time tracking,
		and daily activity reports to ensure accountability and operational transparency.

		\generalItem
		Built a factory maintenance management platform enabling real-time machine status reporting, staged
		maintenance workflows, automated notifications, repair logs, and material usage tracking with automatic
		inventory deduction and procurement insight reporting.

		\generalItem
		Collaborated directly with international stakeholders, conducted iterative demos, incorporated production
		feedback, and deployed finalized systems used by operational teams.
	\end{listBlock}

	\sectionTitle{EDUCATION}

	\subSectionTitle{Gannon University}{Pennsylvania, United States}
	\educationInfo{Bachelor of Science in Computer Systems Engineering}{May 2025\cite{GannonBSDiploma}}

	\subSectionTitle{Harvard University}{Online}
	\educationInfo{CS50’s Introduction to Computer Science}{December 2022\cite{HarvardCS50Certification}}

	\subSectionTitle{Extreme Networks}{Online}
	\educationInfo{Introduction to Future Networks}{May 2022\cite{ExtremeCertification}}

	\sectionTitle{PUBLICATIONS}

	\begin{listBlock}
		\generalItem
		Mossi, R. A. — “Tag AI-Sandbox,” Proceedings of the ASEE Annual Conference, 2025 (peer-reviewed).\cite{AITagResearchPublication}
	\end{listBlock}

    \sectionTitle{TECHNICAL SKILLS}

    \begin{listBlock}

        \generalItem
        Languages: C++, TypeScript, JavaScript, Python, Java, SQL, Bash, PowerShell.

        \generalItem
        Frontend: React, Next.js, Vite, Tailwind CSS, Shadcn/UI.

        \generalItem
        Backend \& Cloud: Supabase, PostgreSQL, AWS S3,
        Amazon Textract, Docker, Auth0, REST APIs.

        \generalItem
        AI \& Computer Vision: OpenAI APIs, Claude APIs, YOLO, CUDA, PyTorch,
        NEAT, Reinforcement Learning, OCR Pipelines, Multimodal AI Workflows.

        \generalItem
        Developer Tooling: GitHub Actions, Husky, ESLint, Prettier,
        Commitizen, CMake, Git.

        \generalItem
        Operating Systems: Linux, Windows, macOS.

        \generalItem
        Languages: Spanish (Native), English (Fluent), Japanese (Intermediate).

    \end{listBlock}

	\sectionTitle{AWARDS AND HONORS}

	\begin{listBlock}
		\awardItem{School Development Award}{award obtained for outstanding development of a videogame using Blueprints in Unreal Engine 4}{2020\cite{SoftwareEngineer2020Award}}

		\awardItem{School Development Award}{award obtained for outstanding development of a videogame using C++ and Unreal Engine 4}{2016\cite{SoftwareEngineer2016Award}}
	\end{listBlock}

	\bibliographyBlock

\end{Form}
\end{document}
